\section{Methodology}
\label{sec:method}

\subsection{Datasets}
\label{sec:datasets}

\para{Primary: \hanna.}
We use the \hanna dataset~\citep{chhun2022hanna}, which contains 1,056 stories generated by 11 systems (10 LLMs from the GPT-2 era plus human-written stories).
Each story is rated by 3 human annotators on 6 dimensions---Relevance, Coherence, Empathy, Surprise, Engagement, and Complexity---using a 1--5 Likert scale.
We define human interestingness as the mean of Surprise and Engagement ratings, averaged across the 3 annotators.
These two dimensions most closely capture the concept of ``interestingness'' as the combination of novelty (Surprise) and sustained attention (Engagement).

\para{Complementary: Chatbot Arena.}
We use a subsample of 200 battles from the Arena Human Preference 55K dataset~\citep{lmsys2023arena}, where users engaged two anonymous chatbots and voted for their preferred response.
This provides a naturalistic pairwise preference signal to complement \hanna's absolute ratings.

\subsection{LLM Judge Evaluations}
\label{sec:judges}

We select three OpenAI models as judges: \gptfo, \gptfomini, and \gptfourmini.
Each judge rates every \hanna story on a 1--5 interestingness scale using a standardized prompt.
We use temperature${}=0$ and seed${}=42$ for reproducibility, yielding $1{,}056 \times 3 = 3{,}168$ story evaluations.
For the \arena subsample, both responses in each battle are rated by all three judges, producing $200 \times 2 \times 3 = 1{,}200$ response evaluations.

\subsection{Metrics}
\label{sec:metrics}

We compute the following per-story metrics:
\begin{itemize}[leftmargin=*,itemsep=0pt,topsep=0pt]
    \item \textbf{LLM consensus}: mean rating across the 3 LLM judges.
    \item \textbf{LLM disagreement}: standard deviation across the 3 LLM judges.
    \item \textbf{Human interestingness}: mean of Surprise and Engagement, averaged across 3 human annotators.
    \item \textbf{Human quality}: mean of all 6 human-rated dimensions (used as a control variable).
\end{itemize}

\subsection{Statistical Tests}
\label{sec:tests}

We test four hypotheses with Bonferroni-corrected significance ($\alpha = 0.0125$):

\para{H1: LLM consensus vs.\ Human Surprise.}
Spearman's $\rho$ between LLM consensus and human Surprise ratings.
The outsider hypothesis predicts a negative correlation.

\para{H2: LLM disagreement vs.\ Human Engagement.}
Spearman's $\rho$ between LLM disagreement (standard deviation) and human Engagement ratings.
The hypothesis predicts a positive correlation---stories that divide LLM judges should engage humans more.

\para{H3: Bottom vs.\ Top LLM consensus quartile.}
Mann-Whitney $U$ test comparing human interestingness ratings between stories in the bottom and top quartiles of LLM consensus.
The hypothesis predicts higher human interest in the bottom quartile.

\para{H4: LLM consensus vs.\ Human interest (LLM-generated only).}
Spearman's $\rho$ restricted to LLM-generated stories only, excluding human-written stories that LLMs might rate low for reasons unrelated to the hypothesis.

\subsection{Quality-Controlled Analyses}
\label{sec:quality_control}

To disentangle quality from interestingness, we compute partial correlations between LLM consensus and human interestingness after regressing out human-rated overall quality (mean of all 6 dimensions).
We also conduct within-model analyses to control for model-level quality differences, and examine extreme groups---stories where all 3 LLM judges agree on low ($\leq 2$) or high ($\geq 4$) ratings.

\subsection{Implementation Details}
\label{sec:implementation}

All experiments use Python 3.12 with NumPy 2.2.5, Pandas 2.2.3, and SciPy 1.15.2.
We use a random seed of 42 throughout.
LLM evaluations are collected via the OpenAI API.
